\section{A second root-bundle derivation of the Kuranishi statement}\label{app:cdp-kuranishi}
We re-derive Theorem~\ref{fam:st:A.5.2}, credited to \citet[Theorem~6.12 and Proposition~6.16]{DuEtAl2026v1}.
\begin{proof}
We confirm the particular implication used here from the root-bundle and cusp-normalization data proved above. On $\mathfrak B=\mathbb P(3,4)$, both vertical and
structure-cohomology bundles are $E=G=\mathcal O\oplus\mathcal O(-1)$.
The vertical subline is generated by the invariant translation
field; its quotient has twelfth power $\mathcal O(-p_0)$ by the
discriminant's simple cusp zero, hence is $\mathcal O(-1)$ in the
torsion-free root Picard group. The extension splits because
$H^1(\mathcal O(1))=0$. For $G$, relative duality and the
normalization complex of the hexagon give
$\det G=R^2\bar f_*\mathcal O_X=\mathcal O(-1)$; the invariant
integral cocycle $\gamma$ is nonzero on every ordinary fibre and
on the dual-graph cycle at the cusp. It gives the same split
extension. This checks the nonvanishing at the cusp required by the structure-sheaf calculation.

The horizontal logarithmic line is $\mathcal O(-5)$. Projection
and Leray applied to
$0\to\bar f^*E\to T_X(-\log W)\to\bar f^*\mathcal O(-5)\to0$
give logarithmic tangent cohomology $(1,1,0,0)$: every nontrivial
root line appearing is one of
$\mathcal O(-1),\mathcal O(-2),\mathcal O(-5),\mathcal O(-6)$,
whose coarse direct image is $\mathcal O_{\mathbb P^1}(-1)$,
while the two surviving trivial summands lie in degrees zero
and one. The branch-normal sequence has quotient
$\eta_*K_{\mathrm dP_6}$, with cohomology $(0,0,1)$; the local
map sends each ambient vector field to its separate normal
coefficient on each branch. Hence its exact cohomology sequence
gives $(1,1,1,0)$ for $T_X$.

The parameter derivative is nonzero. The fixed real marking of
the torus lattice gives the constant Beltrami class
\[
 -\frac{1}{2iD}\,\partial_{\zeta_2}\otimes
 (d\bar\zeta_2-a\,d\bar\zeta_1),
 \quad a=\frac{\operatorname{Im}\mu}{\operatorname{Im}\tau},
 \quad D=\operatorname{Im}\beta_c
       -\frac{6(\operatorname{Im}\mu)^2}{\operatorname{Im}\tau}.
\]
A nonzero constant Dolbeault tensor on a torus is not exact,
because integration of every derivative of a periodic coefficient
is zero. A first-order trivialization of $X_c$ would preserve
the relative complete pencil $|-2K|$. Its induced infinitesimal
automorphism of $\mathbb P^1$ fixes all three critical values:
the equations $t=s^3$, $t-1=s^4$, and $q=xyz$ make each critical
value a constant section over the dual numbers. An infinitesimal
projective transformation fixing three distinct points is zero.
The trivialization would therefore trivialize the nonzero fibre
class just displayed, a contradiction. Thus the period derivative
spans the one-dimensional $H^1(T_X)$.

Write the Kuranishi local algebra as $\mathbb C\{z\}/I$ in its
one-dimensional tangent space. The existing proper period family
defines a local classifying map $z\mapsto\phi(s)$ with
$\phi'(0)\ne0$. The analytic inverse-function theorem makes
$\mathbb C\{z\}\to\mathbb C\{s\}$ an isomorphism. Since it
annihilates $I$, necessarily $I=0$. This verifies smoothness and
reducedness as well as miniversality, and every sufficiently small
deformation is a member of that period family.
\end{proof}